\section*{Hoofdstuk \ref{ch:b3:chromatin-epigenetics} --- Chromatine en epigenetica}

\begin{solution}{exo:b3:chromatin-epigenetics:1}
Van elk van H2A, H2B, H3 en H4 twee kopieën (een octameer) met
\qty{147}{bp} DNA in $1.65$ slagen; linkers van \qtyrange{20}{60}{bp},
wat een herhaling van bijna \qty{200}{bp} geeft. Het nuclease knipt
alleen het blootliggende linker-DNA, nooit binnen een kerndeeltje,
zodat elk fragment een geheel aantal herhalingen is: een ladder van
\num{200}, \num{400}, \qty{600}{bp}, geen veeg.
\end{solution}

\begin{solution}{exo:b3:chromatin-epigenetics:2}
Diploïd: $6.0\times 10^{9}$ bp; $6.0\times 10^{9}/190 \approx
3.2\times 10^{7}$ \omterm{def:b3:chromatin-epigenetics:nucleosome}{nucleosomen}; lengte $6.0\times 10^{9}\times
\qty{0.34}{nm} = \qty{2.0}{m}$.
\end{solution}

\begin{solution}{exo:b3:chromatin-epigenetics:3}
H3K27me3: \omterm{def:b3:chromatin-epigenetics:nucleosome}{histon} H3, lysine 27, getrimethyleerd --- zwijgend
\omterm{def:b3:chromatin-epigenetics:nucleosome}{chromatine} van onderdrukte ontwikkelingsgenen, geschreven door PRC2 en
gelezen door de chromodomein-eiwitten van PRC1. H4K16ac: H4, lysine
16, geacetyleerd --- actief, open \omterm{def:b3:chromatin-epigenetics:nucleosome}{chromatine} (het verbreekt een
contact tussen twee \omterm{def:b3:chromatin-epigenetics:nucleosome}{nucleosomen}). H3S10ph: H3, serine 10,
gefosforyleerd --- een mitotisch merkteken dat samengaat met de
condensatie van chromosomen (en, kortstondig, met genactivatie door
sommige signalen); geen stabiel uitschakelend of activerend merkteken.
\end{solution}

\begin{solution}{exo:b3:chromatin-epigenetics:4}
Het oranjegen is X-gebonden met een oranje en een zwart allel; een
gevlekte vacht vereist van elk één, op twee verschillende
X-chromosomen, waarvan er per cel één zwijgt. Een kater heeft één X en
is helemaal oranje of helemaal zwart. Een mannelijke lapjeskat draagt
twee X-chromosomen en een Y (XXY), en is steriel.
\end{solution}

\begin{solution}{exo:b3:chromatin-epigenetics:5}
$\qty{200}{bp}\times \qty{0.34}{nm} = \qty{68}{nm}$ per kraal van
\qty{11}{nm}: graad $6.2$. Vereist: $\qty{4}{cm}/\qty{4}{\micro m} =
10^{4}$. Verdere vouwing met een factor $10^{4}/6.2 \approx 1600$.
\end{solution}

\begin{solution}{exo:b3:chromatin-epigenetics:6}
$\mu = 0.98$, $\delta = 0.01$: $p^{*} = 0.01/0.03 = 0.33$; verhouding $\mu -
\delta = 0.97$; halveringstijd $\ln 2/(-\ln 0.97) = 22.8 \approx 23$
delingen. $\mu = 0.999$, $\delta = 0.001$: $p^{*} = 0.5$; verhouding
$0.998$; halveringstijd $\ln 2 / 0.002 \approx 350$ delingen.
\end{solution}

\begin{solution}{exo:b3:chromatin-epigenetics:7}
$A$: actief (promotor met H3K4me3 en \omterm{def:b3:chromatin-epigenetics:histone-code}{acetylering}). $B$: zwijgend, door
Polycomb onderdrukt --- facultatief \omterm{def:b3:chromatin-epigenetics:nucleosome}{heterochromatine}, de toestand van
een ontwikkelingsgen dat in deze lijn is uitgezet maar in andere wordt
gebruikt, zodat $B$ het gen is dat in een ander celtype actief kan
zijn. $C$: zwijgend, constitutief \omterm{def:b3:chromatin-epigenetics:nucleosome}{heterochromatine} (de HP1-route),
meestal herhalingen of genen die in elk celtype zwijgen.
\end{solution}

\begin{solution}{exo:b3:chromatin-epigenetics:8}
\emph{UBE3A} komt in neuronen alleen vanaf het moederlijke allel tot
expressie. Haar deletie kwam van haar vader, op het allel dat toch al
zwijgt, dus is zij gezond; geeft zij hem door, dan wordt het een
moederlijke deletie: elk kind heeft kans $1/2$ hem te erven en heeft
dan het syndroom van Angelman. Was de deletie van haar moeder gekomen,
dan zou zij zelf het syndroom van Angelman hebben; een gezonde
draagster kan hem alleen van vaderskant hebben geërfd. (Een vaderlijke
deletie van de hele regio zou Prader--Willi geven; \emph{UBE3A} alleen
is daarvoor niet genoeg.)
\end{solution}

\begin{solution}{exo:b3:chromatin-epigenetics:9}
Een X zonder \emph{\omterm{def:b3:chromatin-epigenetics:x-inactivation}{Xist}} kan niet worden geïnactiveerd, dus zwijgt de
andere X in elke cel: de inactivatie is volledig maar niet langer
willekeurig (scheef), en het embryo is levensvatbaar. Beide X zonder
\emph{\omterm{def:b3:chromatin-epigenetics:x-inactivation}{Xist}}: geen inactivatie, twee actieve X-chromosomen, tweevoudige
X-gebonden expressie, en het vrouwelijke embryo sterft vroeg. Een
\emph{\omterm{def:b3:chromatin-epigenetics:x-inactivation}{Xist}}-transgen op een autosoom bekleedt dat autosoom in cis en
legt zijn genen stil, wat functieverlies van één autosomale kopie
geeft --- het experiment dat liet zien dat \emph{\omterm{def:b3:chromatin-epigenetics:x-inactivation}{Xist}} voldoende is
voor uitschakeling in cis.
\end{solution}

\begin{solution}{exo:b3:chromatin-epigenetics:10}
Gemethyleerde cytosines deamineren tot thymine, dat het herstel niet
herkent, zodat een gemethyleerd \omterm{def:b3:chromatin-epigenetics:methylation}{CpG} met hoge snelheid tot TpG of CpA
muteert. Alleen sequenties die generatie na generatie in de
\emph{kiembaan} ongemethyleerd blijven behouden hun \omterm{def:b3:chromatin-epigenetics:methylation}{CpG}'s; de eilanden
hebben de hunne behouden, dus zijn zij het grootste deel van de
geschiedenis van de gewervelden in de kiembaan ongemethyleerd geweest.
Een gen dat alleen in de lever gemethyleerd is, is in zaadcel en
eicel ongemethyleerd; mutaties die in levercellen ontstaan worden niet
geërfd, dus laat somatische methylering geen spoor in de sequentie na.
\end{solution}

\begin{solution}{exo:b3:chromatin-epigenetics:11}
Bij de replicatie worden de oude, gemerkte \omterm{def:b3:chromatin-epigenetics:nucleosome}{histonen} over de twee
dochterduplexen verdeeld en afgewisseld met nieuwe, ongemerkte, zodat
elk dochterdomein ongeveer half gemerkt is. HP1 dat aan een behouden
gemerkt \omterm{def:b3:chromatin-epigenetics:nucleosome}{nucleosoom} bindt trekt SUV39H aan, dat de ongemerkte buren
methyleert; zolang de gemerkte \omterm{def:b3:chromatin-epigenetics:nucleosome}{nucleosomen} dicht genoeg staan om bij
elk gat een \omterm{def:b3:chromatin-epigenetics:histone-code}{schrijver} te brengen, is het domein vóór de volgende
deling weer gevuld --- dezelfde terugkoppeling als de hoge $\mu$ van
de stelling. De uitbreiding stopt bij grenzen: isolatoreiwitten
(CTCF), sterk getranscribeerde genen en tRNA-genen, \omterm{def:b3:chromatin-epigenetics:nucleosome}{nucleosomen} met
onverenigbare actieve merktekens (H3K4me3, \omterm{def:b3:chromatin-epigenetics:histone-code}{acetylering}) die de
\omterm{def:b3:chromatin-epigenetics:histone-code}{schrijver} niet kan gebruiken, en het punt waar de \omterm{def:b3:chromatin-epigenetics:histone-code}{wissers}
(demethylasen, verloop van HDAC's) de \omterm{def:b3:chromatin-epigenetics:histone-code}{schrijver} voorbijstreven. Toets:
trek HP1 kortstondig naar een reportergen (een fusie met een
DNA-bindend domein waarvan de binding met een geneesmiddel wordt
geregeld), laat het weer los, en volg de stilte van het reportergen en
H3K9me3 over de delingen; verwacht een voortduren dat van de \omterm{def:b3:chromatin-epigenetics:histone-code}{schrijver}
afhangt, en verlies wanneer de \omterm{def:b3:chromatin-epigenetics:histone-code}{schrijver} of de dimerisatie van de
\omterm{def:b3:chromatin-epigenetics:histone-code}{lezer} gemuteerd is, of wanneer een isolator binnen het domein wordt
geplaatst.
\end{solution}

\begin{solution}{exo:b3:chromatin-epigenetics:12}
Sluit uit: een gedeeld milieu (de eigen voeding en armoede van de
kleinkinderen); een baarmoedereffect --- de dochter van de aan de
hongersnood blootgestelde moeder was een foetus met haar kiemcellen al
aanwezig, zodat de kiembaan van het kleinkind rechtstreeks werd
blootgesteld en het effect niet transgenerationeel is; genetische
verschillen tussen families die het overleefden en families die dat
niet deden; omgekeerde oorzakelijkheid (lichaamsmassa verandert zelf
de methylering in bloed); verschillen in de samenstelling van de
bloedcellen; en meervoudig toetsen over vele \omterm{def:b3:chromatin-epigenetics:methylation}{CpG}'s. Muizenexperiment:
stel ingeteelde F$_0$-\emph{mannetjes} bloot (geen route via de
baarmoeder), paar ze met niet-blootgestelde vrouwtjes of gebruik
in-vitrofertilisatie met embryotransfer, en volg F$_2$ en F$_3$ langs
de vaderlijke lijn, met meting van de methylering in de zaadcellen van
elke generatie; een merkteken dat in de zaadcellen en het fenotype van
F$_3$ blijft, in een stam zonder genetische variatie, is overdracht
via de kiembaan.
\end{solution}

\begin{solution}{pb:b3:chromatin-epigenetics:1}
\textbf{1.} $6.4\times 10^{9}\times \qty{0.34}{nm} = \qty{2.2}{m}$;
$6.4\times 10^{9}/200 = 3.2\times 10^{7}$ \omterm{def:b3:chromatin-epigenetics:nucleosome}{nucleosomen}.
\textbf{2.} \omterm{def:b3:chromatin-epigenetics:nucleosome}{Histon}: $3.2\times 10^{7}\times 1.08\times 10^{5} =
3.5\times 10^{12}$ Da $= \qty{5.7}{pg}$. DNA: $6.4\times 10^{9}\times
650 = 4.2\times 10^{12}$ Da $= \qty{6.9}{pg}$. Vergelijkbare massa's.
\textbf{3.} Kern: $\tfrac{4}{3}\pi\,(\qty{5}{\micro m})^{3} =
\qty{520}{\micro m^3}$. Kerndeeltje: $\pi\,(5.5)^{2}\times 5.5 =
\qty{520}{nm^3} = 5.2\times 10^{-7}\,\unit{\micro m^3}$; totaal
$3.2\times 10^{7}\times 5.2\times 10^{-7} = \qty{17}{\micro m^3}$:
ongeveer \qty{3}{\%} van de kern.
\textbf{4.} $3.2\times 10^{7}\times \qty{11}{nm} = \qty{0.35}{m}$;
\omterm{prop:b3:chromatin-epigenetics:compaction}{pakkingsgraad} $2.2/0.35 \approx 6$.
\textbf{5.} $6.4\times 10^{9}/2\times 10^{5} = \num{32000}$ lussen van
elk \num{1000} \omterm{def:b3:chromatin-epigenetics:nucleosome}{nucleosomen}.
\textbf{6.} $0.21^{2}\times 6.4\times 10^{9} = 2.8\times 10^{8}$;
waargenomen $5.6\times 10^{7}$, dus \qty{80}{\%} verloren, door
deaminering van 5-methylcytosine tot thymine (niet-herstelde
C$\to$T-transities op gemethyleerde \omterm{def:b3:chromatin-epigenetics:methylation}{CpG}'s).
\textbf{7.} $p_{n+1} = 0.99\,p_{n} + 0.02\,(1 - p_{n}) = 0.97\,p_{n} +
0.02$; $p^{*} = 0.02/0.03 = 0.67$.
\textbf{8.} $p_{n} = 0.67 + 0.33\times 0.97^{n}$: $n = 10$: $0.91$;
$n = 50$: $0.74$; $n = 200$: $0.67$ (het overschot is $7\times 10^{-4}$).
\textbf{9.} $\ln 2/(-\ln 0.97) = 22.8 \approx 23$ delingen; bij één per
week ongeveer \num{23} weken --- vijf maanden.
\textbf{10.} Elke replicatie koppelt één oude streng, nog gemethyleerd,
aan één nieuwe streng die geen enzym methyleert; er komen geen oude
strengen bij en er gaan er geen verloren, het aantal strengen
verdubbelt, dus halveert de gemethyleerde fractie. $0.80/2^{n} < 0.05$
vereist $2^{n} > 16$: vijf delingen ($0.025$; vier geven precies
\qty{5}{\%}).
\textbf{11.} Verhouding $0.9995 - 0.0005 = 0.999$; halveringstijd $\ln
2/(-\ln 0.999) \approx 690$ delingen, ongeveer \num{13} jaar bij één
per week.
\textbf{12.} Een gen dat een leven lang van zo'n \num{4000} wekelijkse
delingen zwijgt heeft een halveringstijd van honderden delingen nodig,
en die geeft alleen de terugkoppeling van vraag~11: \omterm{def:b3:chromatin-epigenetics:histone-code}{lezers} van het
merkteken (MeCP2, HP1) trekken zijn \omterm{def:b3:chromatin-epigenetics:histone-code}{schrijvers} aan (H3K9-methylase,
DNMT3), zodat een dicht domein van merktekens zichzelf kopieert met
een efficiëntie die geen enkel enzym alleen haalt.
\textbf{13.} De ene X draagt het oranje allel, de andere het zwarte;
elke huidcel heeft één X uitgeschakeld en brengt alleen de andere tot
expressie; de afstammelingen van een cel houden haar keuze, dus is een
plek een kloon (of een groep aangrenzende klonen) die dezelfde keuze
maakte.
\textbf{14.} Alle $N$ kiezen de X met oranje, of alle kiezen de andere:
$2\times (1/2)^{N} = 2^{1-N}$.
\textbf{15.} $2^{1-N} = 10^{-9}$: $N - 1 = \log_{2} 10^{9} = 29.9$,
$N \approx 30$ cellen.
\textbf{16.} $\qty{3000}{cm^2}/\qty{15}{cm^2} = 200$ plekken, veel meer
dan $30$ stichters: tijdens de groei worden de afstammelingen van een
stichter door celmenging en migratie verspreid, zodat één kloon
verscheidene eilanden vormt (en buren met dezelfde keuze versmelten,
wat de andere kant op werkt maar minder).
\textbf{17.} Eén \omterm{def:b3:chromatin-epigenetics:x-inactivation}{Barr-lichaampje} per X boven de eerste: XXY één, XXX
twee, XY geen.
\textbf{18.} Verdragen wanneer het gen een werkzame Y-homoloog heeft,
zodat mannetjes ook twee kopieën hebben (de pseudoautosomale genen en
paren als \emph{KDM5C}/\emph{KDM5D}); het doet ertoe bij aneuploïdieën
van X --- bij het syndroom van Turner (XO) zijn de ontsnapte genen in
enkelvoud aanwezig, en daarom is een XO-individu niet eenvoudig een
normaal vrouwtje (haplo-insufficiëntie van \emph{SHOX} en kleine
gestalte), en bij XXY zijn de ontsnappers overgedoseerd.
\textbf{19.} $1/\num{15000} = 6.7\times 10^{-5}$; deleties $0.70\times
6.7\times 10^{-5} = 4.7\times 10^{-5}$ (één op \num{21000});
moederlijke disomie $1.7\times 10^{-5}$ (één op \num{60000}).
\textbf{20.} Angelman is het verlies van één gen, \emph{UBE3A}, dus
volstaat één puntmutatie in het moederlijke allel. Prader--Willi is het
verlies van verscheidene vaderlijk tot expressie komende genen tegelijk
(\emph{SNRPN}, het cluster kleine RNA's); een puntmutatie schakelt er
één van uit en geeft hoogstens een gedeeltelijk fenotype, nooit het
hele syndroom.
\textbf{21.} \qty{20}{\%} van de kinderen krijgt een vaderlijke deletie
en heeft Prader--Willi; geen enkel kind heeft Angelman, waarvoor een
moederlijke deletie nodig is.
\textbf{22.} Drie chromosomen, twee moederlijke en één vaderlijke,
waarvan er willekeurig één verloren gaat: het vaderlijke gaat verloren
met kans $1/3$, wat moederlijke disomie en Prader--Willi oplevert; met
kans $2/3$ gaat een moederlijke kopie verloren en is het kind normaal.
\textbf{23.} Vaderlijke genen zouden de vraag aan de moeder moeten
opdrijven; hun verlies (Prader--Willi) zou de vraag moeten verlagen ---
wat past bij het slechte zuigen en het niet gedijen van de pasgeborene
--- terwijl de latere onverzadigbare eetlust het tegenovergestelde is
en moeilijker te plaatsen (een voorgestelde lezing is dat de
vaderlijke genen het kind normaal naar het spenen toe sturen).
Angelman, het verlies van een moederlijk gen, zou de vraag moeten
verhogen; het langdurige zuigen en het veelvuldig wakker worden van
Angelman-kinderen passen daarbij, de epilepsie en het verlies van
spraak gaan helemaal niet over hulpbronnen.
\textbf{24.} Neem als doelwit het vaderlijke antisense-transcript
(\emph{UBE3A-ATS}) met een antisense-oligonucleotide dat het afbreekt
of zijn verlenging blokkeert: het vaderlijke \emph{UBE3A} is intact en
wordt alleen in cis door dat RNA uitgeschakeld, dus herstelt het
weghalen van het RNA zijn expressie. De methylering van de
\omterm{def:b3:chromatin-epigenetics:imprinting}{imprintingcontroleregio} blijft onaangeroerd, zodat \emph{SNRPN} en de
andere geïmprinte genen hun normale ouderlijke patroon houden.
\textbf{25.} Halveringstijd van een merkteken: ongeveer $23$ delingen
zonder terugkoppeling, ongeveer $690$ met terugkoppeling; de
lapjesvacht werd beslist door een stichterspopulatie van ongeveer $30$
cellen.
\end{solution}
